\capitulo{7}{Conclusiones y Líneas de trabajo futuras}

\section{Conclusiones}

El presente Trabajo de Fin de Grado ha tenido como objetivo el desarrollo de un sistema capaz de registrar y analizar la interacción del usuario en entornos de modelado tridimensional, utilizando Blender como plataforma de referencia. A lo largo del proyecto se ha diseñado e implementado un addon que permite capturar información relevante del proceso de modelado de forma automática y no intrusiva.

A partir del trabajo realizado, se pueden extraer las siguientes conclusiones principales:

\begin{itemize}

    \item Se ha demostrado la viabilidad de registrar de forma continua la interacción del usuario dentro de Blender, capturando datos relacionados con transformaciones, geometría, eventos de modelado y cambios de estado sin interferir en el flujo de trabajo.

    \item El sistema desarrollado permite estructurar la información en un formato adecuado para su posterior análisis, facilitando la obtención de métricas temporales, espaciales y geométricas del proceso de modelado.

    \item La integración del sistema dentro del propio entorno de Blender, mediante un addon, permite una recopilación de datos en condiciones reales de uso, a diferencia de otros enfoques basados en entornos experimentales o simulados.

    \item Se ha conseguido detectar eventos relevantes de modelado, como duplicados, fusiones de geometría, cambios de modo o modificaciones en coordenadas UV, lo que permite caracterizar de forma más precisa el comportamiento del usuario.

    \item El uso de mecanismos de persistencia de datos integrados en el propio archivo \textit{.blend} garantiza la continuidad de la información entre sesiones, evitando la pérdida de datos y facilitando su gestión.

\end{itemize}

En conjunto, el sistema desarrollado proporciona una base sólida para el análisis del proceso de modelado tridimensional, permitiendo estudiar la evolución de la escena y el comportamiento del usuario a lo largo del tiempo.

\section{Conclusiones técnicas}

Desde un punto de vista técnico, el desarrollo del proyecto ha permitido obtener las siguientes conclusiones:

\begin{itemize}

    \item El uso de la API de Blender (\texttt{bpy}) permite acceder de forma eficiente a las estructuras internas del entorno, facilitando la extracción de información relevante de la escena y de los objetos.

    \item La utilización de \textit{handlers} y temporizadores permite implementar un sistema de captura de eventos en tiempo real, adaptado al funcionamiento interno de Blender.

    \item La detección de cambios mediante comparación de estados (snapshots) resulta una estrategia eficaz para evitar el registro redundante de información y reducir el volumen de datos generados.

    \item La implementación de un sistema de logging no intrusivo requiere un equilibrio entre frecuencia de captura y rendimiento, siendo necesario optimizar el sistema para minimizar el impacto en la experiencia del usuario.

    \item La gestión de eventos complejos, como las modificaciones en coordenadas UV o las operaciones de modelado encadenadas, requiere técnicas específicas de detección y procesamiento que van más allá del simple registro de comandos.

\end{itemize}

Estas conclusiones ponen de manifiesto la complejidad asociada a la captura y estructuración de datos en entornos interactivos, así como la necesidad de diseñar soluciones adaptadas al funcionamiento interno de las herramientas utilizadas.

\section{Limitaciones del sistema}

A pesar de los resultados obtenidos, el sistema desarrollado presenta ciertas limitaciones que deben ser consideradas:

\begin{itemize}

    \item El sistema se ha desarrollado específicamente para Blender, lo que limita su aplicabilidad directa a otros entornos de modelado 3D.

    \item La detección de algunos eventos complejos depende del estado interno del software, lo que puede generar pequeñas imprecisiones en situaciones específicas.

    \item El volumen de datos generado puede ser elevado en sesiones prolongadas, lo que requiere estrategias de optimización y filtrado.

    \item El análisis de los datos recogidos se ha planteado de forma básica, siendo necesario el desarrollo de herramientas más avanzadas para su explotación completa.

\end{itemize}

\section{Líneas de trabajo futuras}

A partir del trabajo realizado, se plantean diversas líneas de mejora y desarrollo futuro:

\begin{itemize}

    \item \textbf{Mejora del sistema de análisis}: desarrollo de herramientas avanzadas para el análisis de los datos recogidos, incluyendo técnicas de minería de datos, aprendizaje automático o visualización interactiva en 3D.

    \item \textbf{Visualización de métricas}: implementación de sistemas que permitan representar los datos directamente dentro del entorno tridimensional, facilitando su interpretación mediante gráficos, mapas o animaciones.

    \item \textbf{Extensión a otros entornos}: adaptación del sistema a otras herramientas de modelado 3D o CAD, permitiendo comparar procesos de diseño entre diferentes plataformas.

    \item \textbf{Optimización del rendimiento}: mejora de la eficiencia del sistema para reducir el impacto en el rendimiento en escenas complejas o sesiones prolongadas.

    \item \textbf{Análisis del comportamiento del usuario}: aplicación del sistema para estudiar patrones de diseño, diferencias entre usuarios o evolución del aprendizaje en contextos educativos.

    \item \textbf{Integración con inteligencia artificial}: utilización de los datos recopilados para entrenar modelos capaces de predecir acciones del usuario, sugerir operaciones o asistir en el proceso de modelado.

\end{itemize}

En conjunto, estas líneas de trabajo abren nuevas posibilidades para el estudio del proceso de diseño y el desarrollo de herramientas inteligentes que mejoren la interacción entre el usuario y los entornos de modelado tridimensional.